% GENERATED FILE. Edit canonical lessons, then run npm run generate:books.
% canonical-source-hash: fnv1a64:2562244207680953
% canonical-lessons: RU-C262-snachala-potom, RU-C262-rasskazat, RU-C262-chtenie, RU-C262-review

\chapter{Then, and At Last: A Story Told in Order}
\label{ch:ru-then-and-at-last-a-story-told-in-order}
\hlchaptermodality{262}

\begin{chapteropening}
\textbf{By the end of this chapter:} \emph{I can tell a short story in order with \ru{сначала}, \ru{потом}, \ru{затем} and \ru{наконец}, retell it with \ru{рассказать}, and read a connected story in the past.}

Complete RU-C262-review: a day retold in four steps, \ru{сначала}, \ru{потом}, \ru{затем}, \ru{наконец}, with the scene in the imperfective and the events in the perfective.
\end{chapteropening}

\section[\texorpdfstring{\ru{сначала}, \ru{потом} … (snachála, potóm …)}{сначала, потом … (snachála, potóm …)}]{\ru{сначала}, \ru{потом}, \ru{затем}, \ru{наконец} — the steps of a story}

\label{lesson:RU-C262-snachala-potom}



\begin{quote}
Say \textquotedblleft{}at first\textquotedblright{}, \textquotedblleft{}later, then\textquotedblright{} and \textquotedblleft{}at last\textquotedblright{}. (\emph{Snachála. Potóm. Nakonéts.}) Three step words you have had for a while. Now they get a story to carry.
\end{quote}

\subsection*{You'll want to know: \ru{затем}}
\textbf{\ru{затем}} — \emph{zatém} — \textquotedblleft{}after that, then\textquotedblright{}. It means much the same as \textbf{\ru{потом}}, and sounds a little more careful: it is common in writing, in instructions and in a story told with some care. Between them, a story need not say \textbf{\ru{потом}} twice in a row.

\begin{grammarlens}[title={Grammar Lens: one finished step, then the next}]
\begin{quote}
\textbf{\ru{Сначала} \ru{я} \ru{позавтракал}, \ru{потом} \ru{прочитал} \ru{газету}.} — \emph{snachála ya pozávtrakal, potóm prochitál gazétu} — First I had breakfast, then I read the paper. \textbf{\ru{Затем} \ru{я} \ru{позвонил} \ru{маме}. \ru{Наконец} \ru{я} \ru{пошёл} \ru{на} \ru{работу}.} — \emph{zatém ya pozvoníl máme. nakonéts ya poshól na rabótu} — After that I rang Mum. At last I set off for work.
\end{quote}

Every verb is perfective, so every step is done before the next one starts. \emph{\ru{Позавтракал}} is \emph{\ru{завтракать}} with \emph{\ru{по}-}. \textbf{\ru{Пошёл}}, \textquotedblleft{}set off\textquotedblright{}, is the past of \emph{\ru{пойти}}, the perfective partner of \emph{\ru{идти}}.
\end{grammarlens}

\begin{culture}
\begin{quote}
\textbf{\ru{Обычно} \ru{я} \ru{сначала} \ru{завтракал}, \ru{потом} \ru{читал} \ru{газету}.} — \emph{obýchno ya snachála závtrakal, potóm chitál gazétu} — Usually I had breakfast first, then read the paper.
\end{quote}

Imperfective verbs and \textbf{\ru{обычно}}: not one morning, but every morning.
\end{culture}

\begin{cousinweb}
\textbf{\ru{затем}} is \emph{\ru{за}}, \textquotedblleft{}behind, after\textquotedblright{}, with \emph{\ru{тем}}, a form of \textbf{\ru{тот}}, \textquotedblleft{}that\textquotedblright{}: \textquotedblleft{}after that\textquotedblright{}. EtymDB 2.1 has no entry for it, so this is the plain reading of its parts. \textbf{\ru{Тот}} sits on the Proto-Indo-European pointing stem \emph{*to-}, the same one behind English \emph{that} and \emph{the}.
\end{cousinweb}

\subsection*{Your turn}
Say these aloud:

\begin{itemize}
\raggedright
  \item \emph{zatém}, stress on the end
  \item \emph{Snachála ya pozávtrakal, potóm prochitál gazétu.}
  \item \emph{Zatém ya pozvoníl máme. Nakonéts ya poshól na rabótu.}
  \item the same four steps as a woman: \emph{pozávtrakala}, \emph{prochitála}, \emph{pozvoníla}, \emph{poshlá}
  \item your own morning in three steps, \emph{snachála}, \emph{potóm}, \emph{nakonéts}
\end{itemize}

\begin{tcolorbox}[breakable,colback=teal!4,colframe=teal!35!black,title={Before you move on}]
Another word for \emph{\ru{потом}}, a little more careful? (\emph{\textbf{\ru{Затем}.}}) One morning, step by step: which aspect? (\textbf{Perfective.}) Every morning, as a habit? (\textbf{Imperfective}, often with \emph{\textbf{\ru{обычно}}}.)
\end{tcolorbox}
\section[\texorpdfstring{\ru{рассказать} (rasskazát)}{рассказать (rasskazát)}]{\ru{рассказать} (rasskazát) — to tell a story, once and to the end}

\label{lesson:RU-C262-rasskazat}



\begin{quote}
Say \textquotedblleft{}to tell (a story)\textquotedblright{}. (\emph{Rasskázyvat.}) Then three steps of a morning, with \emph{snachála}, \emph{potóm} and \emph{zatém}.
\end{quote}

\subsection*{You'll want to know: \ru{рассказать}}
\textbf{\ru{рассказать}} — \emph{rasskazát} — \textquotedblleft{}to tell\textquotedblright{}, once, the whole story. It is the perfective partner of \textbf{\ru{рассказывать}}.

\begin{quote}
\textbf{\ru{Бабушка} \ru{часто} \ru{рассказывала} \ru{сказки}.} — \emph{bábushka chásto rasskázyvala skázki} — Grandma often told fairy tales. \textbf{\ru{Вчера} \ru{бабушка} \ru{рассказала} \ru{новую} \ru{сказку}.} — \emph{vcherá bábushka rasskazála nóvuyu skázku} — Yesterday Grandma told a new fairy tale.
\end{quote}

The habit takes the imperfective; the one tale, told to its end, the perfective. To ask for a story: \textbf{\ru{Расскажи}!} — \emph{rasskazhí!} — Tell me! To a stranger, or to more than one: \textbf{\ru{Расскажите}!} — \emph{rasskazhítye!}

\begin{grammarlens}[title={Grammar Lens: scene, steps, end}]
A told story has three parts, and aspect marks each one.

The \textbf{scene}, in the imperfective: \textbf{\ru{Это} \ru{было} \ru{в} \ru{субботу}. \ru{Шёл} \ru{дождь}.} — \emph{éto býlo v subbótu. shol dozhd} — It was on a Saturday. It was raining. Russian rain \textquotedblleft{}goes\textquotedblright{}: \emph{\ru{идёт} \ru{дождь}}, and in the past \emph{\ru{шёл} \ru{дождь}}.

The \textbf{steps}, in the perfective and in order, with \emph{\ru{сначала}}, \emph{\ru{потом}}, \emph{\ru{затем}}. Then the \textbf{end}, with \textbf{\ru{наконец}}.
\end{grammarlens}

\begin{cousinweb}
\textbf{\ru{рассказать}} is \emph{\ru{рас}-}, \textquotedblleft{}out, apart\textquotedblright{}, with \textbf{\ru{сказать}}, \textquotedblleft{}to say\textquotedblright{}. By the usual account \emph{\ru{сказать}} is built on an older \emph{\ru{казать}}, \textquotedblleft{}to show\textquotedblright{}, which EtymDB 2.1 traces to Proto-Slavic \emph{*kazati}; so telling is showing in words. EtymDB lists \emph{\ru{рассказать}} itself without taking it apart.
\end{cousinweb}

\subsection*{Your turn}
Say these aloud:

\begin{itemize}
\raggedright
  \item \emph{rasskazát}, then the past: \emph{rasskazál}, \emph{rasskazála}
  \item \emph{Vcherá bábushka rasskazála nóvuyu skázku.}
  \item \emph{Rasskazhí!}, then to a stranger: \emph{Rasskazhítye!}
  \item a scene: \emph{Éto býlo v subbótu. Shol dozhd.}
  \item three steps after it: \emph{snachála …, potóm …, nakonéts …}
\end{itemize}

\begin{tcolorbox}[breakable,colback=teal!4,colframe=teal!35!black,title={Before you move on}]
\textquotedblleft{}Grandma told a fairy tale\textquotedblright{}, once? (\emph{\textbf{\ru{Бабушка} \ru{рассказала} \ru{сказку}.}}) Every evening? (\emph{\textbf{\ru{Рассказывала}.}}) Which aspect sets the scene? (\textbf{The imperfective.}) And the steps? (\textbf{The perfective.})
\end{tcolorbox}
\section[(reading)]{(reading) — a rainy Saturday, told in order}

\label{lesson:RU-C262-chtenie}



\begin{quote}
The book's first passage was a few lines in the present. This one is a story, and it happened last summer and last Saturday.

Every word in it has come up before. Three verbs appear in a new shape, and they are explained after the passage.
\end{quote}

\begin{tcolorbox}[breakable,skin=enhanced,colback=orange!6,colframe=orange!45!black,boxrule=0.5pt,arc=1mm,left=6pt,right=6pt,top=4pt,bottom=4pt,fonttitle=\bfseries,title={Read this}]
\begin{quote}
\ru{Это} \ru{было} \ru{в} \ru{субботу}. \ru{Шёл} \ru{дождь}. \ru{Я} \ru{сидел} \ru{в} \ru{кафе}, \ru{пил} \ru{кофе} \ru{и} \ru{читал} \ru{газету}. \ru{Когда} \ru{я} \ru{прочитал} \ru{газету}, \ru{пришла} \ru{моя} \ru{подруга} \ru{Анна}. \ru{Сначала} \ru{мы} \ru{пили} \ru{чай} \ru{и} \ru{говорили}. \ru{Потом} \ru{Анна} \ru{рассказала} \ru{историю}.
\end{quote}

\begin{quote}
\ru{Летом} \ru{она} \ru{была} \ru{в} \ru{Москве} \ru{и} \ru{жила} \ru{в} \ru{гостинице} \ru{в} \ru{центре}. \ru{Утром} \ru{она} \ru{часто} \ru{гуляла} \ru{в} \ru{парке}. \ru{Но} \ru{в} \ru{понедельник}, \ru{пока} \ru{она} \ru{гуляла}, \ru{она} \ru{потеряла} \ru{паспорт}! \ru{Сначала} \ru{она} \ru{плакала}. \ru{Потом} \ru{она} \ru{позвонила} \ru{в} \ru{гостиницу}. \ru{Затем} \ru{она} \ru{вернулась} \ru{в} \ru{парк}. \ru{Наконец} \ru{она} \ru{нашла} \ru{паспорт}. \ru{Он} \ru{был} \ru{в} \ru{сумке}!
\end{quote}

\begin{quote}
\ru{Мы} \ru{смеялись}. \ru{Вечером} \ru{мы} \ru{ужинали} \ru{в} \ru{ресторане}, \ru{и} \ru{Анна} \ru{рассказала} \ru{ещё} \ru{одну} \ru{историю}.
\end{quote}

Read it once without stopping. Do not translate as you go.

Then read it again and listen for the step words.
\end{tcolorbox}

\begin{culture}
The scene is imperfective: \emph{\ru{сидел}}, \emph{\ru{пил}}, \emph{\ru{читал}}, \emph{\ru{шёл} \ru{дождь}}. The steps are perfective, in order: \emph{\ru{прочитал}}, \emph{\ru{пришла}}, \emph{\ru{рассказала}}, \emph{\ru{потеряла}}, \emph{\ru{позвонила}}, \emph{\ru{вернулась}}, \emph{\ru{нашла}}.

Three of them are new shapes of verbs you own. \textbf{\ru{Пришла}} is the perfective partner of \emph{\ru{приходить}}, \textquotedblleft{}she came\textquotedblright{}. \textbf{\ru{Вернулась}} is the partner of \emph{\ru{возвращаться}}, \textquotedblleft{}she went back\textquotedblright{}. \textbf{\ru{Нашла}} is the partner of \emph{\ru{находить}}, \textquotedblleft{}she found\textquotedblright{}.

\emph{\ru{Часто} \ru{гуляла}} is a habit, every morning. \emph{\ru{Пока} \ru{она} \ru{гуляла}} is one walk going on; \emph{\ru{потеряла}} is the moment it went wrong. \emph{\ru{Плакала}} and \emph{\ru{смеялись}} are imperfective: crying and laughing that went on for a while, with no finish to report.
\end{culture}

\subsection*{Your turn}
\begin{itemize}
\raggedright
  \item \emph{Say it:} the passage aloud, once, without stopping
\end{itemize}

\begin{tcolorbox}[breakable,colback=teal!4,colframe=teal!35!black,title={Before you move on}]
How many step words keep Anna's day in order? (\textbf{Four:} \emph{\textbf{\ru{сначала}, \ru{потом}, \ru{затем}, \ru{наконец}}}.)

\emph{\ru{Пока} \ru{она} \ru{гуляла}, \ru{она} \ru{потеряла} \ru{паспорт}}: which verb was going on? (\emph{\textbf{\ru{Гуляла}.}}) Which one happened once? (\emph{\textbf{\ru{Потеряла}.}})
\end{tcolorbox}
\section[(review)]{Review — a day, retold in order}

\label{lesson:RU-C262-review}



\begin{quote}
Say the four step words in order. (\emph{Snachála, potóm, zatém, nakonéts.})
\end{quote}

\begin{grammarlens}[title={What you've built}]
Three lessons, and a whole story fits in your mouth. A scene in the imperfective: \emph{\ru{было}}, \emph{\ru{шёл} \ru{дождь}}, \emph{\ru{сидел}}. Steps in the perfective, one finished before the next, with \textbf{\ru{сначала}}, \textbf{\ru{потом}}, \textbf{\ru{затем}} and \textbf{\ru{наконец}}. And \textbf{\ru{рассказать}} for telling it once, to the end.

Anna's day, in four steps:

\begin{quote}
\textbf{\ru{Сначала} \ru{она} \ru{плакала}. \ru{Потом} \ru{она} \ru{позвонила} \ru{в} \ru{гостиницу}. \ru{Затем} \ru{она} \ru{вернулась} \ru{в} \ru{парк}. \ru{Наконец} \ru{она} \ru{нашла} \ru{паспорт}.}
\end{quote}

The crying went on, so it is imperfective. The call, the walk back and the finding were each done once.
\end{grammarlens}

\subsection*{The exchange}
\begin{quote}
\textbf{— \ru{Расскажи}, \ru{что} \ru{было} \ru{в} \ru{субботу}!} — \emph{rasskazhí, shto býlo v subbótu!} — Tell me what happened on Saturday! \textbf{— \ru{Мы} \ru{были} \ru{в} \ru{кафе}. \ru{Сначала} \ru{мы} \ru{пили} \ru{чай}, \ru{потом} …} — \emph{my býli v kafé. snachála my píli chay, potóm …}
\end{quote}

\subsection*{Your turn}
\begin{itemize}
\raggedright
  \item \emph{Say it:} Anna's four steps, from memory
  \item \emph{Say it:} \emph{Rasskazhí, shto býlo v subbótu!}
  \item \emph{Answer:} your own Saturday: one scene, then three steps with \emph{snachála}, \emph{potóm}, \emph{nakonéts}
  \item \emph{Say it:} \emph{Kogdá ya chitál, pozvoníla Ánna.} One going on, one cutting in
\end{itemize}

\begin{tcolorbox}[breakable,colback=teal!4,colframe=teal!35!black,title={Before you move on}]
The step word between \emph{\ru{потом}} and \emph{\ru{наконец}}? (\emph{\textbf{\ru{Затем}.}}) \textquotedblleft{}Tell me!\textquotedblright{}? (\emph{\textbf{\ru{Расскажи}!}}) The scene, or the steps: which takes the perfective? (\textbf{The steps.})
\end{tcolorbox}
